\documentclass[10pt,a4paper]{amsart}
\usepackage[margin=19mm]{geometry}
\usepackage{amsmath,amssymb,mathtools}
\usepackage[scr=rsfs,cal=euler]{mathalfa}
\usepackage{xfrac}
\usepackage[unicode,hidelinks,bookmarks=false]{hyperref}
\newcommand{\ie}{i.e.\ }
\newcommand{\cf}{cf.\ }
\newcommand{\resp}{respectively\ }
\newcommand{\Subsection}{\subsection}
\providecommand{\nobreakdash}{\nobreak\hbox{-}}
\setlength{\emergencystretch}{2em}
\multlinegap0pt
\usepackage{bm}
\usepackage{bbm}
\usepackage{dsfont}
\usepackage{xspace}
\usepackage{cancel}
\usepackage{mathtools}
\usepackage{amsthm}
\makeatletter
\@ifundefined{theo}{\newtheorem{theo}{Theo}}{}
\@ifundefined{lem}{\newtheorem{lem}[theo]{Lemma}}{}
\makeatother
\title[The Maker-Breaker directed triangle game]{The Maker-Breaker directed triangle game}
\author{Hrishikesh Jagtap, Moumanti Podder}
\date{Source: arXiv:2510.13919v2 (CC BY 4.0)}
\begin{document}
\maketitle

\noindent\emph{Source and licence.} The Maker-Breaker directed triangle game. Version arXiv:2510.13919v2 (CC BY 4.0), distributed under the Creative Commons Attribution 4.0 International licence, as recorded in the arXiv licence metadata for this exact version and shown on the abstract page https://arxiv.org/abs/2510.13919v2. Full rights notice: licenses/arxiv-2510.13919-CC-BY-4.0.txt.\par\smallskip

\noindent\textit{Editorial scope.} Counted unit: the Lem whose source label is \texttt{lem:existence} in the article (source0.tex lines 896--902, source label \texttt{lem:existence}), delivered together with the article's own complete original proof of it (source0.tex lines 904--1008, one \texttt{proof} environment of 105 lines). No mathematics is written, repaired or re-derived: statement and proof are the article's own text line for line. Required context retained uncounted: none. Every definition, display and label of those lines is retained unchanged. Omitted from this package and not redistributed: 1--895, 903--903, 1009--1428. Nothing inside a retained excerpt is omitted or altered: every retained body is delivered exactly as the article's lines are written, so no omission receipt of any kind accompanies this package. Citations: the retained excerpts carry no citation command at all, so no bibliography entry of the article is transcribed and no reference is redistributed. Reference adaptation: none was needed and none was made; every reference inside the delivered unit resolves inside it, so no unresolved reference or citation is produced. Printed slips kept literal: no printed word, symbol, hypothesis or number of the counted statement or of its proof was altered. Layout and class: the article's own class is \texttt{article} with packages including PRIMEarxiv, inputenc, fontenc, fullpage, graphicx, amsmath, amssymb, amsfonts, amsthm, mathtools, booktabs, nicefrac, microtype, lipsum; the wrapper is the lane's pinned \texttt{amsart} editorial wrapper at 10pt on A4, which restates the theorem-like environments the retained text needs and adds the article's own macro definitions that the retained text needs (none), with extra wrapper packages (bm, bbm, dsfont, xspace, cancel, mathtools). The delivered numbering is the unit's own. Assets: no retained span contains a figure environment, a graphics-inclusion command, a PDF-page inclusion, a table or a TikZ picture; no image file is copied into this package, the package declares no asset dependency and no image file is redistributed with it. Licence: version arXiv:2510.13919v2 of this article is distributed under the Creative Commons Attribution 4.0 International licence, as recorded in the arXiv licence metadata for this exact version and shown on the abstract page https://arxiv.org/abs/2510.13919v2. A full rights notice is delivered at licenses/arxiv-2510.13919-CC-BY-4.0.txt. Characters: every retained excerpt is pure ASCII, so the delivered engine, XeLaTeX with its default Latin Modern text font, reports no missing character.
\begin{lem}\label{lem:existence}
Let $T(n,p)$ be a random tournament with fixed $p\in(0,1)$.
Then
\[
\lim_{n\to\infty}\mathbb P\bigl(T(n,p)\text{ contains an induced copy of }\Pi(7)\bigr)=1.
\]
\end{lem}

\begin{proof}
For each $7$-element subset $A=\{v_1,\dots,v_7\}\subset V_n$ with
\[
v_1<v_2<\cdots<v_7,
\]
let $I_A$ be the indicator random variable of the event that the tournament induced by $T(n,p)$ on $A$
is exactly $\Pi(7)$ under the order-preserving identification $v_i\mapsto i$ for $i=1,\dots,7$.

Let
\[
X_n=\sum_{\substack{A\subset V_n\\ |A|=7}} I_A,
\]
so that $X_n$ counts the number of induced copies of $\Pi(7)$ in $T(n,p)$.

Among the pairs $1\le i<j\le 7$, exactly $12$ satisfy $i+j$ odd and exactly $9$ satisfy $i+j$ even.
Therefore, for every $A$,
\[
\mathbb P(I_A=1)=p^{12}(1-p)^9=:q>0.
\]
Hence
\[
\mathbb E[X_n]=\binom{n}{7}q=\Theta(n^7).
\]

We now estimate the second moment.
Write
\[
\mathbb E[X_n^2]
=
\sum_{\substack{A,B\subset V_n\\ |A|=|B|=7}}
\mathbb E[I_A I_B].
\]
For $0\le r\le 7$, let
\[
S_r
=
\sum_{\substack{A,B\subset V_n\\ |A|=|B|=7\\ |A\cap B|=r}}
\mathbb E[I_A I_B].
\]
Then
\[
\mathbb E[X_n^2]=\sum_{r=0}^7 S_r.
\]

If $r=0$, then $A$ and $B$ are disjoint, so the relevant edge orientations are independent. Thus
\[
\mathbb E[I_A I_B]=\mathbb E[I_A]\mathbb E[I_B]=q^2,
\]
and therefore
\[
S_0=\binom{n}{7}\binom{n-7}{7}q^2.
\]
Since
\[
\frac{\binom{n-7}{7}}{\binom{n}{7}}\to 1,
\]
it follows that
\[
S_0=(1+o(1))\binom{n}{7}^2 q^2=(1+o(1))(\mathbb E[X_n])^2.
\]

Now fix $1\le r\le 7$.
The number of ordered pairs $(A,B)$ of $7$-subsets with $|A\cap B|=r$ is
\[
\binom{n}{r}\binom{n-r}{7-r}\binom{n-7}{7-r}=O(n^{14-r}).
\]
Also,
\[
0\le \mathbb E[I_A I_B]\le 1.
\]
Hence
\[
S_r=O(n^{14-r})=o(n^{14})=o\bigl((\mathbb E[X_n])^2\bigr),
\]
because $(\mathbb E[X_n])^2=\Theta(n^{14})$.

Combining the estimates for $S_0,\dots,S_7$, we obtain
\[
\mathbb E[X_n^2]=(1+o(1))(\mathbb E[X_n])^2.
\]
Therefore
\[
\operatorname{Var}(X_n)
=
\mathbb E[X_n^2]-(\mathbb E[X_n])^2
=
o\bigl((\mathbb E[X_n])^2\bigr).
\]

By Chebyshev's inequality,
\[
\mathbb P(X_n=0)
\le
\mathbb P\bigl(|X_n-\mathbb E[X_n]|\ge \mathbb E[X_n]\bigr)
\le
\frac{\operatorname{Var}(X_n)}{(\mathbb E[X_n])^2}
\longrightarrow 0
\qquad\text{as }n\to\infty.
\]
Thus
\[
\mathbb P(X_n>0)\longrightarrow 1,
\]
which means that $T(n,p)$ contains an induced copy of $\Pi(7)$ with probability tending to $1$.
\end{proof}

\end{document}
